\scanpage{184}
\mspage{58}
% archive sheet number: 183

\[
M^{\unsure{\pi|\Pi}}=M\cap\mathrm{Filt}^{\Pi}.)
\]

\inserted{c$'$) Il existe une partie $M_0$ de $\mathrm{Filt}^{\Pi}$, telle que $M$ soit la plus petite partie de $\mathcal M$ stable par limites ind. filtrantes. (De plus, on peut alors prendre $M_0$ \uline{petit})}

% slanted
\margin{\uline{NB} On
\diagram{\begin{tikzcd}[ampersand replacement=\&,column sep=tiny,row sep=small]
 a\arrow[r,Leftrightarrow]\arrow[d,Rightarrow] \& b\arrow[d,Rightarrow] \& \deleted{\unsure{\Rightarrow c}}\\
 c')\arrow[r,Rightarrow,"{\textcircled{?}}"] \& c) \& \deleted{\unsure{\Rightarrow}}
\end{tikzcd}}
\unsure{sans} condition sur $\pi$ (sauf $\pi\geq\pi_0$)}

\uline{Corollaire 5.8.} Les applications $M\longmapsto M^\pi=M\cap\mathrm{Filt}^{\Pi}$,

$M^\pi\longmapsto$ l'ensemble des $X\in\mathrm{Ob}\,\mathcal M$ tels que l'ens $\mathrm{Ss}^{M^\pi}_{\pi}(X)$ des $Z_\alpha\in\mathrm{Ss}_{\pi}(X)$ qui sont dans $M^\pi$ soit cofinal
\inserted{($=$ \deleted{\unsure{le}} \deleted{plus} petit sous-ens. de $\mathrm{Ob}\,\mathcal M$ stable par $\varinjlim$ filtrantes)}

\unsure{établissent|et définissent} des bijections, inverses l'une de l'autre, entre l'ens des $M\subset\mathrm{Ob}\,\mathcal M$ satisfaisant les conditions \inserted{équivalentes} de la prop. 5.7, et l'ensemble des parties $M^\pi$ de $\mathrm{Filt}^{\Pi}$ stables par $\varinjlim$ filtrantes de cardinal $\leq\pi$.

\uline{Dém. de 5.7.} L'implication b) $\Rightarrow$ c) est tautologique, compte tenu de la définition de la $\pi$-accessibilité de $M$, et de la prop. 5.3. \deleted{L'implication c) $\Rightarrow$ b) n'est autre que prop. \unsure{5.5.} (\unsure{revue|vue} et \unsure{corrigée}), que je viens de prouver. Donc $b\Leftrightarrow c$,} Et $a\Rightarrow b$ provient de la stabilité de $\mathrm{Filt}^{\pi}$ par $\underline{\lim}$ filtrantes de cardinal $\leq\pi$ (condition d)).

Il reste donc à prouver b) $\Rightarrow$ a) \inserted{(C'est ici qu'on doit utiliser l'hyp. particulière sur $\pi$)} \deleted{ce qui se réduit à prouver que b) implique que $M$ est stable par (petites) $\varinjlim$ filtrantes.} Une fois ceci vu, on saura l'équivalence de a,b,c, et \deleted{d'autre part cela implique aussitôt (via a et b)} d'autre part il est immédiat que a) $\Rightarrow$ c$'$), en

